\begin{table}[htbp]
  \centering
  \caption{Flächenanteil je Brennstoffmodell oberhalb einsatzrelevanter Flammenlängen, in Prozent.}
  \label{tab:schwellen}
  \begin{tabular}{lrrrrrrrr}
    \toprule
    Brennstoffmodell & \multicolumn{4}{c}{$>1{,}2$\,m} & \multicolumn{4}{c}{$>3{,}4$\,m} \\
    \cmidrule(lr){2-5}\cmidrule(lr){6-9}
     & S1 & S2 & S3 & S4 & S1 & S2 & S3 & S4 \\
    \midrule
    Heide (GA-heath) & 0{,}0 & 96{,}9 & 100{,}0 & 100{,}0 & 0{,}0 & 0{,}2 & 96{,}1 & 96{,}7 \\
    Gras (GR4) & 0{,}0 & 98{,}5 & 99{,}8 & 100{,}0 & 0{,}0 & 0{,}0 & 0{,}3 & 0{,}9 \\
    Laubwald (TL6) & 0{,}0 & 0{,}1 & 0{,}1 & 0{,}1 & 0{,}0 & 0{,}0 & 0{,}0 & 0{,}0 \\
    Nadelwald (TU5) & 5{,}2 & 17{,}3 & 86{,}7 & 97{,}7 & 0{,}9 & 2{,}4 & 8{,}1 & 37{,}1 \\
    \bottomrule
  \end{tabular}
  \par\smallskip\footnotesize Schwellenwerte nach \textcite{andrewsrothermel1982}: oberhalb von 1{,}2\,m (4\,ft) ist ein direkter Angriff mit Handkräften nicht mehr möglich, oberhalb von 3{,}4\,m (11\,ft) sind Kontrollversuche an der Feuerfront in der Regel wirkungslos.
\end{table}
